\section{A bounded table for an unbounded suffix}
\label{sec:bounded-suffix}

A fixed table horizon does not by itself suffice for an unbounded request
stream: a long interval may contain many requests while the optimum hardly
moves. We will retain a bounded word for the table and restart it only after
enough offline movement has accrued to absorb a fixed per-epoch cost.
Throughout this section the finite table is assumed available;
Section~\ref{sec:uniform-computation} schedules its construction.

The routine may begin at any deterministic request index, from arbitrary
actual server positions.  Partition the suffix into \emph{raw blocks}:
a block ends at the request that first raises its number of distinct
requested points to $k+1$, and the next request starts a new block.
Every completed block forces any service to pay at least $\delta$, since
its entering $k$ server positions cannot cover all $k+1$ requested points.
Set
\begin{equation}\label{eq:epoch-length}
 R=\left\lceil\frac{kD}{\delta}\right\rceil
\end{equation}
and let an \emph{epoch} consist of $R$ completed raw blocks. A completed
epoch therefore forces cost at least $R\delta\ge kD$; this lower bound
will absorb restart charges proportional to $kD$. A finite input may end inside a
block or at a block boundary; every nonempty epoch then contains at most
$R$ nonempty blocks, counting a final incomplete block. Block and epoch
boundaries depend only on raw requests and can be recognized before
serving their last request.

Within each epoch, we follow the finite table until its number of positive
moves exceeds a fixed cap. After the cap is exceeded, a marking rule will
serve the remainder at a cost of at most $(k+1)D$ per block
(Lemma~\ref{lem:fallback}). Accordingly, choose
\begin{equation}\label{eq:suffix-parameters}
 F=R(k+1)D,\qquad
 M=\left\lceil\frac{F}{\delta}\right\rceil,\qquad
 H=\sum_{i=0}^{M}k^i.
\end{equation}
Here $F$ reserves the entire fallback cost, and $M+1$ positive table moves
cost more than $F$. The filter below records capped trajectories in a tree
with at most $H$ nodes and charges each retained request to a different
node. Use the dyadic table of Proposition~\ref{prop:finite-table} for this
$H$, and write $a$ for its coefficient. All parameters are fixed by the
instance, independently of the request sequence and its length.

\subsection{A universal capped filter}

The filter, cap, and fallback mechanism is related to the request filtering
and reset construction of Mikkelsen~\cite[Section~5.2,
Lemmas~8--9]{Mikkelsen2016}. We prove the particular branching-tree version
needed here, including a bound on all created entries. That bound will
control both the table horizon and the storage needed during service.

At the start of each epoch initialize a list with the single entry
$(u,0)$.  An entry $(c,i)$ consists of a distinct-position labeled tuple
and a nonnegative move count.  Separate entries with identical data are
permitted.  On each raw request $r$ apply these deterministic rules:
\begin{enumerate}[label=(\roman*),leftmargin=*]
\item If every listed tuple covers $r$, drop the request and leave the
list unchanged.  This rule also applies to an empty list.
\item Otherwise keep the request.  Preserve entries that cover $r$.
Remove every other entry $(c,i)$; if $i<M$, replace it by the $k$ entries
$(T(c,r,j),i+1)$, for $j=1,\ldots,k$.  An entry with $i=M$ is removed
without replacement.
\end{enumerate}
Maintain the word of kept requests.  The filter is run on every raw
request of the epoch, including requests served by the fallback below;
it uses no random bits and does not inspect the virtual or actual service.

\begin{lemma}[Universal filter]\label{lem:filter}
After any raw prefix, the list represents exactly the endpoint tuples
and move counts of all strongly lazy services from $u$ of the kept word
using at most $M$ positive moves.  Over the entire epoch the filter creates
at most $H$ entries, counting removed entries.  In particular, its live
list and its kept word both have size at most $H$, regardless of the
number of raw requests.
\end{lemma}

\begin{proof}
The representation assertion follows by induction.  Dropping a request
does not alter the kept word.  On a kept request, a strongly lazy service
at a covering tuple makes no move and retains its tuple and count.  At a
noncovering tuple it must choose one of the $k$ labels, increasing the
positive-move count by one.  These are exactly the replacements in the
rule, with paths exceeding the cap discarded.  Starting from distinct
positions ensures that every resulting tuple remains distinct.

Give the initial entry a root node.  When an entry is removed, call its
node spent; if replacements are made, give them $k$ distinct child nodes.
The resulting rooted tree has branching at most $k$ and depth at most
$M$, so it has at most $\sum_{i=0}^{M}k^i=H$ nodes.  A hit keeps the same
node, and a spent node never returns.  Every kept request spends at least
one previously unspent node.  This proves the bound on the kept word as
well as the bound on all created entries.  It includes the last request
that might empty the list; all subsequent requests are then dropped.
\end{proof}

The quantifier over \emph{all} strongly lazy services is useful here.
In particular, the list includes every legal rounded-table trajectory
within the move cap, even one that visits a configuration having zero
probability before the table's dyadic rounding.

\subsection{Virtual service and fallback}

At every epoch start, initialize a virtual tuple at $u$, a positive-move
counter at zero, and the filter and kept prefix as above.  Virtual service
has two modes, initially table mode.  Every virtual choice of a label is
also the actual algorithm's choice: it moves that actual labeled server
to the current request and leaves the others fixed.  Actual positions
are never reset.

In \emph{table mode}, feed each kept request to the table, using its
pre-request kept prefix and current virtual tuple and fresh unbiased
bits.  On a dropped request, choose the label virtually occupying the
request.  Count positive virtual moves.  If serving a request brings
this count to $M+1$, enter fallback mode \emph{after} serving that
request, for the remainder of the epoch.

These instructions are valid on every random branch.  Before a
table-mode request, the virtual tuple is the endpoint of the table path
on the kept prefix and its move count is at most $M$.  By
Lemma~\ref{lem:filter} it occurs in the pre-request list.  A dropped
request is therefore covered and does not change this tuple.  A kept
request continues the table path legally, possibly making move $M+1$
and triggering the switch.  Lemma~\ref{lem:filter} also guarantees that
no table decision beyond its horizon $H$ is requested.  In particular,
an empty list cannot cause an uncovered dropped request in table mode.

In \emph{fallback mode}, serve raw requests by marking points. This is
the protection rule of paging marking algorithms
\cite[Section~3]{FKLMSY1991}; we use its deterministic bound on the number
of moves within a block. Begin
with an empty mark set on entry into the remainder of the current raw
block, and again at each subsequent raw-block start.  Each mark has a
different assigned label, kept at that marked point until the block
ends.  Use the following rules, always resolving arbitrary choices by
the smallest label:
\begin{enumerate}[label=(\roman*),leftmargin=*]
\item At a raw-block endpoint, choose a covering virtual label if one
exists, and otherwise choose any virtual label.  Serve the request and
discard the marks.
\item At any other request, if the point is marked, choose its assigned
label.  If it is unmarked, choose a covering virtual label if possible;
otherwise choose a label not assigned to a mark.  Serve the request,
mark its point, and assign the chosen label to it.
\end{enumerate}
An already covering label at a new point is unassigned, since labels
assigned to previous marks remain at those different points.  All
virtual rules preserve distinct positions.

The raw-block set and completed-block counter are updated using
\emph{every} raw request, independently of modes and the filter.  If
the switch occurs at a block endpoint, fallback marks start empty in
the next block, if any.  If it occurs at an epoch endpoint, there is
no fallback remainder in that epoch.  In all cases, the next epoch
begins in table mode with a fresh virtual tuple and fresh filter;
these are changes to internal state only.  This specifies a causal
algorithm without needing to know when the input ends.

\begin{lemma}[Fallback cost]\label{lem:fallback}
The marking rules are always feasible.  Following any switch, their
total virtual cost over the rest of the epoch is at most $F$, including
when the input ends inside a raw block.
\end{lemma}

\begin{proof}
Consider a nonendpoint request at a newly marked point.  Through this
request the whole raw block contains at most $k$ distinct points.  Since
the existing marks are different previously requested points of that
block, there are fewer than $k$ marks before the new point is marked.
An unassigned label therefore exists.  No movement at a nonendpoint
disturbs a previously assigned label, so repeated marked requests cost
zero.

There are at most $k$ distinct nonendpoint points in a raw block.  Each
costs at most $D$ on its first fallback appearance, and the endpoint,
if served in fallback, costs at most one further $D$.  Thus any remaining
block portion costs at most $(k+1)D$.  This argument also covers a raw
endpoint reached before its fallback portion has seen $k+1$ distinct
points: the endpoint rule simply discards its marks early.  An incomplete
final block has no endpoint cost.  At most $R$ block portions remain in
the epoch, proving the bound $R(k+1)D=F$.  The request that triggers the
switch belongs entirely to table mode and is not counted here.
\end{proof}

\subsection{Comparisons and amortization}

The filter and fallback now have bounds independent of the raw epoch
length. To combine them, we compare the virtual service with a complete
table trajectory on a deterministic word. This avoids conditioning the
table estimate on the random switch time. We then pay separately for the
actual entering configuration and for each epoch restart.

Fix a finite request sequence $\sigma$ independently of all random bits,
and fix the deterministic index at which the routine begins.  Choose
one optimal labeled service of the \emph{whole} sequence from $s$.
For each nonempty epoch $E$, including its possibly truncated final
portion, let $O_E$ be this optimum's movement during $E$, and let $w_E$
be the filter's full kept word in $E$.  Both the epoch boundaries and
$w_E$ are deterministic.

\begin{lemma}[Epoch comparison]\label{lem:epoch-comparison}
For every such epoch, from any actual entering tuple, the expected actual
cost of the routine satisfies
\begin{equation}\label{eq:epoch-cost}
 \E[\operatorname{cost}_{\rm actual}(E)]
 \le 2aO_E+(2a+2)kD.
\end{equation}
The bound also holds if the entering tuple depends on earlier random bits.
\end{lemma}

\begin{proof}
Run a shadow copy of the table from $u$ on all of $w_E$, and let $V_E$
be its cost.  More precisely, assign independent blocks of $b$ unbiased
bits to the successive positions of this deterministic word.  Use the
same blocks for the actual table-mode decisions; blocks after a switch
are needed only for the shadow's analysis.  Dropped requests leave the
table-mode tuple unchanged, so the shadow and virtual service agree
through every table-mode decision, including a triggering move.
The shadow has the ordinary full-table law, without conditioning on
whether or when the switch occurs; Figure~\ref{fig:epoch-shadow} depicts
this comparison. Proposition~\ref{prop:finite-table}
and Lemma~\ref{lem:filter} give
\begin{equation}\label{eq:shadow-table}
 \E[V_E]\le a\operatorname{OPT}_u(w_E)+D.
\end{equation}

Table-mode cost is at most $V_E$.  If there is a switch, its common
prefix already contains $M+1$ positive moves, each costing at least
$\delta$, and hence
\[
 V_E\ge (M+1)\delta\ge F.
\]
Lemma~\ref{lem:fallback} then pays for the entire fallback by a second
copy of $V_E$.  If there is no switch, that extra copy is unnecessary.
Thus, pathwise,
\begin{equation}\label{eq:virtual-shadow}
 \operatorname{cost}_{\rm virtual}(E)\le 2V_E.
\end{equation}
This includes a switch on the final raw request or final kept request,
and any truncation of an epoch.  The expectation in
Equation~\eqref{eq:shadow-table} is always applied to a fixed full word.

For each label, its first use in the epoch brings both its actual and
virtual positions to the same requested point.  By the triangle
inequality its actual cost on that use exceeds its virtual cost by at
most $D$, the maximum distance between their entering positions.  All
later costs for that label agree, until the next virtual reset.  Unused
labels cost nothing.  Therefore
\begin{equation}\label{eq:actual-virtual}
 \operatorname{cost}_{\rm actual}(E)
 \le \operatorname{cost}_{\rm virtual}(E)+kD.
\end{equation}
Equation~\eqref{eq:actual-virtual} also covers a virtual hit on which the actual chosen label must
move, and it allows repeated actual positions.

We next compare the table's optimum with the fixed whole-sequence
optimum.  For each label $j$, keep only its visits at times retained in
$w_E$, and shortcut its other visits within $E$.  Starting these retained
visits from its original epoch-entering point costs no more than its
movement during $E$.  Replacing that point by $u_j$ adds at most $D$
on the first retained visit, by the triangle inequality; if there is no
retained visit, the label costs zero.  These $k$ paths constitute a legal
labeled service of $w_E$ from $u$, because at each kept time they use the
label chosen there by the whole-sequence optimum.  Consequently
\begin{equation}\label{eq:subsequence-comparison}
 \operatorname{OPT}_u(w_E)\le O_E+kD.
\end{equation}
Combining Equations~\eqref{eq:shadow-table}--\eqref{eq:subsequence-comparison}
gives
\[
 \E[\operatorname{cost}_{\rm actual}(E)]
 \le 2a(O_E+kD)+2D+kD
 \le 2aO_E+(2a+2)kD,
\]
where the last inequality uses $k\ge2$.  All comparisons involving the
entering actual tuple were pathwise and uniform in that tuple; fresh
table bits give the same shadow law after any earlier history.
\end{proof}

\begin{figure}[tbp]
\centering
\begin{tikzpicture}[
  font=\small,
  box/.style={draw=black!65,rounded corners=2pt,align=center,
    inner sep=6pt,text width=3.05cm,minimum height=1.0cm},
  flow/.style={-{Latex[length=2mm]},thick},
  comparison/.style={flow,dashed,draw=blue!55!black}
]
\node[box,fill=black!3] (raw) at (0,0) {Fixed raw epoch $E$};
\node[box,fill=black!3] (filter) at (4.4,0)
  {Deterministic filter\\runs throughout $E$};
\node[box,fill=black!3] (word) at (8.8,0)
  {Full kept word $w_E$\\$|w_E|\le H$};
\node[box,fill=blue!5] (table) at (4.4,-2.0)
  {Virtual table mode\\prefix of $w_E$};
\node[box,fill=blue!5] (shadow) at (8.8,-2.0)
  {Full table shadow\\on $w_E$; cost $V_E$};
\node[box,fill=orange!8] (fallback) at (4.4,-4.0)
  {Marking fallback\\remaining cost $\le F$};
\node[align=center,text width=3.2cm] at (8.8,-4.0)
  {A switch implies\\$V_E\ge(M+1)\delta\ge F$.};
\draw[flow] (raw) -- (filter);
\draw[flow] (filter) -- (word);
\draw[flow] (filter) -- node[left,font=\footnotesize]{kept requests} (table);
\draw[flow] (word) -- (shadow);
\draw[comparison] (shadow) -- node[above=5mm,align=center,font=\footnotesize]
  {same bits\\before switch} (table);
\draw[flow] (table) -- node[right,align=left,font=\footnotesize]
  {after positive\\move $M+1$} (fallback);
\draw[flow] (raw.south) -- (0,-4.0) -- node[above,align=center,font=\footnotesize]
  {raw requests\\after switch} (fallback.west);
\end{tikzpicture}
\caption{The filter is independent of all sampled choices. The full shadow
is an analysis device on the deterministic word $w_E$; it continues after
the virtual service switches to fallback. Shared random bits identify the
table-mode prefix with a prefix of that shadow. Hence table-mode cost is at
most $V_E$, and a switch supplies another $V_E$ to pay for fallback.
The diagram shows data and comparison relations, not a common time scale.}
\label{fig:epoch-shadow}
\end{figure}


\begin{proposition}[Unbounded suffix guarantee]\label{prop:suffix}
For every fixed finite sequence $\sigma$, the routine starting at any
deterministic index, from arbitrary actual positions, has expected
suffix cost at most
\begin{equation}\label{eq:suffix-bound}
 (4a+2)\operatorname{OPT}_s(\sigma)+(2a+2)kD.
\end{equation}
Its table horizon and all parameters in
Equations~\eqref{eq:epoch-length}--\eqref{eq:suffix-parameters}
are fixed by the instance.
\end{proposition}

\begin{proof}
As observed when defining the blocks, the fixed whole-sequence optimum
pays at least $\delta$ inside each completed block, regardless of its
entering tuple. The completed blocks are disjoint intervals of request
times.

Let $N$ be the number of nonempty epochs.  The empty-suffix case is
immediate.  Otherwise the first $N-1$ epochs each contain $R$ completed
blocks, so
\[
 (N-1)R\delta\le \operatorname{OPT}_s(\sigma),
 \qquad
 N\le 1+\frac{\operatorname{OPT}_s(\sigma)}{kD},
\]
where $R\delta\ge kD$ follows from
Equation~\eqref{eq:epoch-length}.  Also
$\sum_E O_E\le\operatorname{OPT}_s(\sigma)$ because the epochs are
disjoint parts of the original sequence.  Summing
Equation~\eqref{eq:epoch-cost} therefore yields
\begin{align*}
 \E[\operatorname{cost}_{\rm suffix}]
 &\le 2a\operatorname{OPT}_s(\sigma)+(2a+2)kD\,N\\
 &\le (4a+2)\operatorname{OPT}_s(\sigma)+(2a+2)kD.
\end{align*}
\end{proof}

Only one epoch's fixed charge remains in this bound.  A final unfinished
epoch may contain arbitrarily many raw requests: its kept word is still
bounded by $H$, and its fallback cost is still bounded by $F$.  Thus no
table error or restart charge accumulates with the raw request count.
